\section{Expected Outcomes}

The project is expected to deliver:

\begin{itemize}[leftmargin=1.4em,itemsep=0.25em]
\item a deployed, gamified web-security training platform hosting an intentionally
      vulnerable web application covering OWASP Top~10 vulnerability
      families, with challenges verified by server-side state change (O1),
\item a gamification module (points, badges, levels, achievements, and a
      peer-banded leaderboard) whose design is explicitly justified against the
      documented negative effects of gamification (O2),
\item an automated scoring mechanism that evaluates each attempt from completion,
      accuracy, and time, with every attempt logged as a per-learner performance
      record (O3),
\item an adaptive support layer: a hint-and-resource ladder, plus a simple,
      interpretable machine-learning model that predicts who is likely to struggle
      and recommends a suitably difficult next challenge, so assignments are
      customised to the individual (O4),
\item an instructor dashboard showing per-learner and cohort progress, challenge
      completion, and overall performance, with at-risk learners flagged (O5),
\item an evaluation, from a pilot with a student cohort, of the platform's effect on
      learner engagement and practical cybersecurity skill (O6).
\end{itemize}

\noindent
Taken together, these deliver a training environment that does three things existing
platforms do not do at once: it adapts to each learner, it gamifies responsibly, and
it gives the instructor an early, ML-informed view of who needs help.
